%% Elevark fra mattelab.no – laget automatisk av bygg.py. Rediger gjerne til eget bruk.
%% Kompiler med lualatex.
\documentclass[11pt,a4paper]{article}
\usepackage[norsk]{babel}
\usepackage[a4paper,margin=2.2cm,top=2cm,bottom=2.4cm]{geometry}
\usepackage{fontspec}
\usepackage{unicode-math}
\setmainfont{TeX Gyre Pagella}
\setsansfont{TeX Gyre Heros}[Scale=MatchLowercase]
\setmathfont{TeX Gyre Pagella Math}
\usepackage{xcolor}
\definecolor{blekk}{HTML}{16213A}
\definecolor{dempet}{HTML}{566077}
\definecolor{rute}{HTML}{CFC9B8}
\definecolor{m1}{HTML}{2F4AB0}
\definecolor{m2}{HTML}{99266D}
\usepackage[most]{tcolorbox}
\usepackage{qrcode}
\usepackage{needspace}
\usepackage{fancyhdr}
\usepackage[hidelinks]{hyperref}
\pdfvariable suppressoptionalinfo \numexpr 1+2+4+8+16+32+64+128+256+512\relax
\setlength{\parindent}{0pt}
\setlength{\parskip}{0.4em}
\pagestyle{fancy}
\fancyhf{}
\renewcommand{\headrulewidth}{0pt}
\fancyfoot[L]{\sffamily\footnotesize\color{dempet}mattelab.no · Potensreglene · elevark uten fasit}
\fancyfoot[R]{\sffamily\footnotesize\color{dempet}Side \thepage}

\newcommand{\fase}[2]{\tcbox[on line, colframe=#2, colback=white, coltext=#2, boxrule=0.6pt, arc=2pt,
  boxsep=0pt, left=3pt, right=3pt, top=1.5pt, bottom=1.5pt, fontupper=\sffamily\bfseries\footnotesize]{#1}}
\newcommand{\steg}[4]{%
  \par\noindent{\color{rute}\rule{\linewidth}{0.4pt}}\par\nopagebreak
  \noindent\begin{minipage}[t]{2.1cm}\vspace{0pt}\fase{#1}{#2}\end{minipage}%
  \begin{minipage}[t]{\dimexpr\linewidth-2.1cm\relax}\vspace{0pt}#3\end{minipage}\par
  \vspace{#4}}
\newcommand{\oppgave}[2]{\par\Needspace{12\baselineskip}\vspace{0.6em}%
  {\fontsize{14}{17}\selectfont\bfseries{\color{m1}#1}\enspace #2}\par\vspace{0.2em}}

\begin{document}
\begin{minipage}[t]{0.72\linewidth}\vspace{0pt}
{\sffamily\small\color{dempet}Elevark · 1T · mattelab.no}\par\vspace{0.3em}
{\fontsize{24}{28}\selectfont\bfseries Potensreglene}\par\vspace{1em}
{\sffamily\small Navn: \rule{4.6cm}{0.4pt}\quad Klasse: \rule{1.5cm}{0.4pt}\quad Dato: \rule{1.9cm}{0.4pt}}
\end{minipage}\hfill
\begin{minipage}[t]{0.22\linewidth}\vspace{0pt}\raggedleft
\qrcode[height=2.3cm]{https://mattelab.no/utforskere/potenser.html}\par\vspace{0.2em}
{\sffamily\scriptsize\color{dempet}Utforskeren}
\end{minipage}

\vspace{1em}
{\sffamily\small\color{dempet}Utforskeren ligger på \url{https://mattelab.no/utforskere/potenser.html}}\par
Hver oppgave har tre steg. \textbf{Gjett} først, uten å bruke utforskeren. \textbf{Test} så i utforskeren. \textbf{Forklar} til slutt med egne ord hvorfor det blir slik.

\oppgave{1}{Gange potenser}
\steg{Gjett}{m1}{Skriv $2^3 \cdot 2^4$ som én potens. Mange svarer $2^{12}$. Hva tror du?}{2.2cm}
\steg{Test}{m2}{Vis $2^3 \cdot 2^4$. Tell boksene.}{2.6cm}
\steg{Forklar}{blekk}{Forklar hvorfor vi legger sammen eksponentene og ikke ganger dem. Kontroller med $8 \cdot 16$.}{3.4cm}

\oppgave{2}{Dele potenser}
\steg{Gjett}{m1}{Hva er $\frac{2^5}{2^2}$? Hva er $\frac{2^2}{2^5}$?}{2.2cm}
\steg{Test}{m2}{Vis $\frac{2^5}{2^2}$ og $\frac{2^2}{2^5}$. Hvor mange bokser strykes? Hvor havner de som er igjen?}{2.6cm}
\steg{Forklar}{blekk}{Forklar hvorfor $\frac{2^2}{2^5} = 2^{-3} = \frac{1}{2^3}$.}{3.4cm}

\oppgave{3}{Hvorfor er $a^0 = 1$?}
\steg{Gjett}{m1}{Hva er $\frac{2^3}{2^3}$? Hva sier regelen for divisjon at det skal være?}{2.2cm}
\steg{Test}{m2}{Vis $\frac{2^3}{2^3}$. Prøv også med andre grunntall.}{2.6cm}
\steg{Forklar}{blekk}{Forklar hvorfor $a^0 = 1$ og ikke $0$.}{3.4cm}

\oppgave{4}{Potens av en potens}
\steg{Gjett}{m1}{Er $(2^3)^2$ lik $2^5$ eller $2^6$?}{2.2cm}
\steg{Test}{m2}{Vis $(2^3)^2$. Hvor mange grupper er det, og hvor mange bokser i hver?}{2.6cm}
\steg{Forklar}{blekk}{Forklar forskjellen på $(2^3)^2$ og $2^3 \cdot 2^2$.}{3.4cm}

\oppgave{5}{Mønsteret nedover}
\steg{Gjett}{m1}{$2^3 = 8$, $2^2 = 4$ og $2^1 = 2$. Fortsett mønsteret nedover. Hva blir $2^0$, $2^{-1}$ og $2^{-2}$?}{2.2cm}
\steg{Test}{m2}{Vis mønsteret for grunntallet $2$. Bytt så til grunntallet $10$.}{2.6cm}
\steg{Forklar}{blekk}{Hva gjør du for å komme fra én rad til den neste? Forklar hvorfor $10^{-1} = 0{,}1$.}{3.4cm}

\end{document}
